\subsubsection{Endpoint płatności}
\label{api:payment}

\texttt{POST /payment/pay} przyjmuje obiekt \hyperref[schema:payment]{Payment}
 w ciele JSON i zwraca obiekt \texttt{PaymentResult}.
Przykład płatności BLIK dla istniejącej rezerwacji o saldzie 10,00 PLN:
\begin{verbatim}
{
  "id": "payment-example-1",
  "reservationId": "existing-reservation-id",
  "amount": "10.00",
  "currency": "PLN",
  "paymentDetails": {
    "paymentMethod": "BLIK",
    "blikPayment": {
      "blikCode": "111111",
      "email": "buyer@example.org"
    }
  },
  "paymentTimestamp": "2026-01-01T00:00:00.123Z"
}
\end{verbatim}
Identyfikator rezerwacji w przykładzie należy zastąpić rzeczywistym.
Kontroler nie generuje identyfikatora płatności ani znacznika czasu;
przekazuje wartości dostarczone przez klienta do dalszego przetwarzania.

Odpowiedź po poprawnym zakończeniu:
\begin{verbatim}
{
  "status": "SUCCESSFUL",
  "statusMessage": null
}
\end{verbatim}
\begin{itemize}
  \item \texttt{SUCCESSFUL}: strategia zaakceptowała płatność, zapisano ją
    i zmieniono status rezerwacji na \texttt{PAID}.
  \item \texttt{DENIED}: strategia odmówiła płatności, np. z powodu kodu BLIK
    lub nieuprawnionego agenta.
  \item \texttt{FAILED}: błąd kwoty lub waluty, operatora albo finalizacji;
    szczegóły znajdują się w \texttt{statusMessage}.
\end{itemize}
Są to statusy wyniku biznesowego, a nie kody HTTP.
Błędy walidacji struktury żądania i brak rezerwacji są zgłaszane wyjątkami,
nie jako \texttt{PaymentResult}.

\subsubsection{Walidacja i wybór strategii}

\texttt{PaymentValidator} sprawdza m.in. niepuste pola rezerwacji, kwoty
 i waluty, nieujemną kwotę dziesiętną, obecność czasu oraz dokładnie jeden
obiekt szczegółów zgodny z \texttt{paymentMethod}.
Kod BLIK musi zawierać sześć cyfr, a e-mail przejść walidację formatu.
Dla gotówki wymagany jest niepusty kod agenta; dla karty m.in. numer,
kod bezpieczeństwa złożony z trzech lub czterech cyfr, miesiąc od 1 do 12,
rok większy od 1900 i nazwa właściciela.

\texttt{PaymentAmountValidator} wymaga jednej waluty rezerwacji oraz
kwoty równej jej wyliczonemu saldu. Rezerwacje bez pozycji
 i wielowalutowe nie są obsługiwane.
\texttt{Money} normalizuje kwoty do dwóch miejsc dziesiętnych
 z zaokrągleniem \texttt{HALF\_UP}; zerowe saldo po płatności jest
sprawdzane numerycznie przez \texttt{signum()}, niezależnie od skali
\texttt{BigDecimal}.
Obecna implementacja sumuje ceny zapisane w migawkach ofert:
nie mnoży ich przez \texttt{OrderItem.quantity} i nie odejmuje wcześniejszych
wpłat. Nie jest to obsługa płatności częściowych.

\texttt{PaymentFacade} filtruje wszystkie strategie według
\texttt{isApplicableFor}. Dokładnie jedna musi pasować:
brak dopasowania oznacza nieobsługiwaną płatność, a wiele dopasowań
--- błąd konfiguracji. Zwracana jest jedyna strategia z listy dopasowań,
nie pierwsza strategia z listy wszystkich komponentów.

\subsubsection{Finalizacja, kompensacja i ograniczenia}

Po zaakceptowaniu płatności przez strategię fasada wywołuje
\texttt{PaymentFinalizer}. Metoda transakcyjna najpierw zapisuje
płatność, a następnie zmienia rezerwację z \texttt{RESERVED} na
\texttt{PAID}. Błąd zmiany statusu powoduje wycofanie zapisu płatności.
Szczegóły wybranej metody są zapisywane jako JSON bez opakowania
\texttt{PaymentDetails}; czas jest przechowywany jako liczba milisekund
od epoki. Wyszukiwanie po rezerwacji porządkuje wyniki po czasie i ID.

Operacja strategii poprzedza lokalną transakcję bazodanową.
Jeśli finalizacja zawiedzie, fasada próbuje wykonać
\texttt{cancelPayment}. Wynik \texttt{FAILED} zawiera przyczynę
finalizacji, a przy nieudanej kompensacji także przyczynę anulowania.
Nie jest to transakcja rozproszona gwarantująca zwrot środków.

\begin{itemize}
  \item BLIK jest symulatorem: \texttt{111111} oznacza sukces,
    \texttt{222222} brak autoryzacji w czasie, \texttt{333333} wygaśnięcie,
    \texttt{444444} niedostępność operatora; inne kody oznaczają odmowę.
  \item Gotówka wymaga jednego sprzedawcy w rezerwacji. Kod agenta musi
    zaczynać się od nazwy sprzedawcy. Jest to uproszczona reguła,
    nie rzeczywisty mechanizm uwierzytelniania.
  \item Obsługa kart używa lokalnego symulatora bramki.
    Anulowanie karty nie jest zaimplementowane i zgłasza
    \texttt{UnsupportedOperationException}.
  \item Anulowanie BLIK i gotówki zwraca sukces bez komunikacji
    z zewnętrznym operatorem.
  \item Repozytorium serializuje szczegóły karty, w tym kod bezpieczeństwa.
    To ograniczenie demonstracyjnej implementacji: nie należy używać
    rzeczywistych danych kart ani traktować jej jako gotowej do produkcji.
\end{itemize}

\subsubsection{Testy płatności}

\texttt{PaymentFacadeTest} obejmuje dziesięć testów jednostkowych:
kolejność operacji, błędy walidacji i operatora, odmowę, kompensację
oraz brak, wielokrotność i wybór dopasowanej strategii.
\texttt{PaymentRepositoryTest} sprawdza m.in. zapis i odczyt
trzech metod płatności, kolejność wyników i klucz obcy rezerwacji.

\texttt{PaymentIntegrationTest} uruchamia rzeczywiste komponenty
Spring Boot i osobną bazę SQLite. Pięć scenariuszy obejmuje:
\begin{enumerate}
  \item poprawną płatność BLIK od kontrolera do bazy,
  \item odrzucenie innej waluty bez zmiany danych,
  \item odrzucenie niepoprawnego kodu BLIK,
  \item wspólne zatwierdzenie płatności i statusu rezerwacji,
  \item wycofanie zapisu płatności przy niedozwolonej zmianie statusu.
\end{enumerate}
Testy integracyjne wywołują bezpośrednio komponent kontrolera,
nie sprawdzają routingu HTTP ani serializacji żądań i odpowiedzi.
Testy transakcji nie mają transakcji otaczającej cały test, dzięki czemu
sprawdzają granicę transakcyjną rzeczywistego finalizatora.

Uruchomienie testów płatności:
\begin{verbatim}
./gradlew test --tests 'pl.edu.uj.student.wzorceprojektowe.payment.*'
\end{verbatim}
Budowanie dokumentacji: \texttt{./gradlew buildDocs}.
Dodanie testów do repozytorium nie stanowi potwierdzenia ich wykonania.
